%%%%%%%%%%%%%%%%%%%%%%%%%%%%%%%%%%%%%%%%%
% Cies Resume/CV
% LaTeX Template
% Version 1.1 (20/7/14)
%
% This template has been downloaded from:
% http://www.LaTeXTemplates.com
%
% Original author:
% Cies Breijs (cies@kde.nl)
% https://github.com/cies/resume with extensive modifications by:
% Vel (vel@latextemplates.com)
%
% License:
% CC BY-NC-SA 3.0 (http://creativecommons.org/licenses/by-nc-sa/3.0/)
%
%%%%%%%%%%%%%%%%%%%%%%%%%%%%%%%%%%%%%%%%%

%----------------------------------------------------------------------------------------
%	PACKAGES AND OTHER DOCUMENT CONFIGURATIONS
%----------------------------------------------------------------------------------------

\documentclass[10pt,a4paper]{article} % Font size (10-12pt) and paper size (a4paper, letterpaper, legalpaper, etc)

\include{structure} % Include structure.tex which contains packages and document layout definitions

\hyphenation{Some-long-word} % Specify custom hyphenation points in words with dashes where you would like hyphenation to occur, or alternatively, don't put any dashes in a word to stop hyphenation altogether

\begin{document} 

%----------------------------------------------------------------------------------------
%	NAME AND CONTACT INFORMATION
%----------------------------------------------------------------------------------------

\maintitle{Gabriel Ramirez}{April 9, 1983}  % Your name and date of birth or subtitle

\noindent\href{mailto:luga.ramirez@gmail.com}{luga.ramirez@gmail.com}
\ifdefined\nophone\else\bull \textsmaller{+}31 626-957064\fi
\bull \href{https://www.linkedin.com/in/lugaramirez/}{lugaramirez} \textit{(LinkedIn)}
\bull \href{https://github.com/lugaramirez}{lugaramirez} \textit{(GitHub)}
\bull Marius Bauerstraat 30-B4
\bull Amsterdam, 1062AR
\bull The Netherlands % Your address

\spacedhrule{0.9em}{-0.4em} % Horizontal rule - the first bracket is whitespace before and the second is after

%----------------------------------------------------------------------------------------
%	SUMMARY SECTION
%----------------------------------------------------------------------------------------

\roottitle{Summary} % Root section title

I specialize in simplifying architectures that have grown too complex for the teams that depend on them --- diagnosing where the weight is, making the structural call, and seeing it through. Over 20 years I've done this across healthcare, e-commerce, fintech, banking, and public infrastructure, in Paraguay, the US, and the Netherlands. My background starts in electrical and embedded engineering, which trained a different kind of systems instinct: I think in constraints, failure modes, and physical limits before I think in abstractions. When a system starts collapsing under growth, the root cause is almost always a physical constraint the architecture stopped respecting --- spotting that gap, e.g., latency, memory pressure, is what the first half of my career trained me to see.

\spacedhrule{0.5em}{-0.4em} % Horizontal rule - the first bracket is whitespace before and the second is after

%----------------------------------------------------------------------------------------
%	EXPERIENCE SECTION
%----------------------------------------------------------------------------------------

\roottitle{Experience} % Top level section

\headedsection % Employer name which can include a hyperlink and location/URL on the right side of the page
{\href{https://bol.com}{bol.}}
{\textsc{Utrecht, The Netherlands}} {

\headedsubsection % Job title entry for the current employer
{Software Engineer}
{Nov '23 -- Present}
{\bodytext{
\textit{The core problems: a tightly-coupled legacy monolith blocking cross-team independence, an inherited culture that treated legacy stability as a reason to avoid tests, and an emerging need to govern \acr{AI}-assisted development responsibly at scale.}
\begin{itemize}
  \item Migrated login handling out of the Java monolith into a new Go service; chose an event-based design over direct service coupling, enabling other teams to add post-login behaviour autonomously without coordination overhead.
  \item Demonstrated test-gated automatic releases on a legacy Kotlin service; the approach spread organically, leading to an invitation to join the 4-person company-wide quality community.
  \item Owns quality and security constraints for \acr{AI}-assisted code generation (Claude) company-wide; introduced promptfoo as a pre-release gate for Claude Code skills before company-wide distribution.
  \item Maintains systems at scale (20M+ users, 95k+ req/s) across Go, Java, Kotlin, and Python --- edge router, GraphQL connectors, and the monolith-to-microfrontend migration.
\end{itemize}
}}

}

%------------------------------------------------

\headedsection % Employer name which can include a hyperlink and location/URL on the right side of the page
{\href{https://quin.md/nl}{Quin}}
{\textsc{Amsterdam, The Netherlands}} {

\headedsubsection % Job title entry for the current employer
{Tech Lead}
{Sep '21 -- Oct '23}
{\bodytext{
\textit{The core problems: a microservices architecture collapsing under its own weight --- a feature as simple as showing a doctor their patient list took 6 to 8 months to ship, infrastructure costs exceeded \$2M annually, deployments were quarterly, and onboarding a new developer took nearly a year.}
\begin{itemize}
  \item Diagnosed the shared commons library as the most immediate drag --- accumulated RestTemplate clients for every Kotlin service in the landscape, loaded by every microservice regardless of need; stripped it down, cutting individual service footprint by 20\%+.
  \item Standardized Git workflows, code review practices, and \acr{CI/CD} pipelines across all backend teams.
  \item Bet on a modular monolith over the prevailing microservices orthodoxy; consolidated 70+ Kotlin services into a single event-driven architecture with 5 bounded domains and 3 anti-corruption layers, applying Domain-Driven Design throughout.
  \item The consolidation cut infrastructure costs from \$2M+ to under \$300K (85\%+ savings), accelerated deployments from quarterly to every merge, and reduced onboarding from nearly a year to days. The company, which had entered bankruptcy, recovered and resumed growth.
\end{itemize}
}}

}

%------------------------------------------------

\headedsection % Employer name which can include a hyperlink and location/URL on the right side of the page
{\href{http://backbase.com/}{Backbase}}
{\textsc{Amsterdam, The Netherlands}} {

\headedsubsection
{Tech Lead}
{May '20 -- Aug '21}
{\bodytext{
\textit{The core problems: an integration strategy built around one-off scenario forms with no reusable components across them, and a release process so fragile that every deployment broke something.}
\begin{itemize}
  \item Diagnosed that the scenario forms shared underlying Java integration components (identity verification, \acr{SMS}/email, document signing) being reimplemented repeatedly; redesigned the approach around composable connectors --- personally built 5, the template grew to 12 with the team, and enabled 20+ connectors after departure.
  \item Built the Angular frontend components for the connector layer personally --- identity verification flows, multi-step onboarding forms, and \acr{SMS}/email journeys --- so each connector was a deployable end-to-end banking journey, not just a backend integration point.
  \item Diagnosed the release failures as a testing gap; introduced acceptance tests and hands-on testing practices across the team, cutting re-releases from twice weekly to once per quarter across 5 banking clients.
\end{itemize}
}}

\headedsubsection
{Backend Trainer}
{Jan '20 -- May '20}
{\bodytext{Delivered 3 week-long training sessions on Backbase's core banking platform to approximately 30 client engineers before the role pivoted due to the pandemic.}}

}

%------------------------------------------------

\headedsection % Employer name which can include a hyperlink and location/URL on the right side of the page
{\href{http://swnat.com/}{Software Natura}}
{\textsc{Asuncion, Paraguay}} {

\headedsubsection
{Team Lead}
{Feb '17 -- Dec '19}
{\bodytext{
\textit{The core problems: a Java call center application that couldn't communicate with a \acr{COBOL} mainframe (one month, declared impossible by everyone else), followed by a request to ``update'' the call center that warranted a full architectural diagnosis.}
\begin{itemize}
  \item Built a C-based Linux kernel module to bridge the Java/Struts application to the \acr{COBOL} mainframe --- the only viable path across a hard process boundary --- and delivered within the one-month deadline.
  \item When asked to update the call center, diagnosed it as beyond patching; made the call for a full revamp in Spring Boot and Angular with domain-driven boundaries, and brought in a data specialist for Apache Spark search rather than attempting it without domain expertise.
  \item Diagnosed the Paraguay-side team as lacking process and standards; introduced structured workflows and code quality practices that enabled a 7-person team to collaborate effectively with 5 US-based developers across time zones.
\end{itemize}
}}

\headedsubsection
{\acr{QA} Engineer}
{May '16 -- Jan '17}
{\bodytext{Diagnosed a 100k+ line legacy \acr{CRM} monolith --- serving American community colleges in a 50-developer team --- as having zero test coverage, meaning bugs only surfaced through weeks of end-user feedback; grew coverage from zero to 60\%, shifting detection to build time.}}

}

%------------------------------------------------

\headedsection % Employer name which can include a hyperlink and location/URL on the right side of the page
{Santco \acr{SRL}}
{\textsc{Asuncion, Paraguay}} {

\headedsubsection
{Software Engineer \& Team Lead}
{Nov '14 -- Mar '16}
{\bodytext{
\textit{The core problem: clients with zero tolerance for downtime and no existing software for their most critical operations --- delivery reliability, not technical capability, was the actual constraint.}
\begin{itemize}
  \item Structured a team of up to 8 developers around Agile iteration cycles, replacing ad-hoc delivery with predictable, testable releases.
  \item Delivered 4 process-automation applications on CakePHP and Google Cloud for Paraguay's leading cattle auction house, largest logistics company, and top leather manufacturer --- the auction house trusted the software enough to increase bidding events from weekly to every two days.
  \item Each delivery required learning the client's domain from the ground up before writing a line of code --- bidding logic, logistics routing, production cycles. That domain fluency is what made the automation trustworthy rather than just functional.
\end{itemize}
}}
}

%------------------------------------------------

\headedsection
{Early Career}
{\textsc{Paraguay, 2005--2012}} {
{\bodytext{Before moving full-time into software engineering, I spent roughly seven years building physical systems at national scale: the network and data center infrastructure for Paraguay's 911-equivalent emergency services (a \$500M+ resilient optical network spanning major metropolitan areas, supporting tens of thousands of street cameras and real-time coordination of police, fire, and ambulance dispatch); a multi-tenant virtual machine provisioning platform built from first principles before \acr{AWS} existed, consolidating infrastructure for 15 clients into a shared on-demand model; and network deployments across 500+ banking sites. Running alongside that: the CompuBus research project at Catholic University --- embedded C and C++ on microcontrollers with Linux kernel-level communication modules, delivering a working prototype deployed with a real bus company in Asuncion. Thinking in hardware constraints before software abstractions is not background flavour --- it is the lens.}}
}

%------------------------------------------------

\spacedhrule{-0.2em}{-0.4em} % Horizontal rule - the first bracket is whitespace before and the second is after

%----------------------------------------------------------------------------------------
%	EDUCATION SECTION
%----------------------------------------------------------------------------------------

\roottitle{Education} % Top level section

\headedsection % Employer name which can include a hyperlink and location/URL on the right side of the page
{University of Leuven}
{\textsc{Leuven, Belgium}} {

\headedsubsection % Job title entry for the current employer
{Masters degree in Nanoscience and Nanotechnology}
{2012 -- 2014}
{\bodytext{Area of interest: Bionanotechnology.}}
}

%------------------------------------------------

\headedsection % Employer name which can include a hyperlink and location/URL on the right side of the page
{Catholic University of Asuncion}
{\textsc{Asuncion, Paraguay}} {

\headedsubsection % Job title entry for the current employer
{Electrical Engineering \textnormal{(discontinued)}}
{2007 -- 2010} {}

\headedsubsection % Job title entry for the current employer
{Bachelors in Electrical Engineering}
{2002 -- 2006} {}
}

\spacedhrule{0.5em}{-0.4em} % Horizontal rule - the first bracket is whitespace before and the second is after

%----------------------------------------------------------------------------------------
%	SKILLS SECTION
%----------------------------------------------------------------------------------------

\roottitle{Skills} % Top level section

\inlineheadsection
{Languages:}
{Java, Kotlin, Go, Python, C, C++, JavaScript, \acr{PHP}, \acr{SQL}.}

\inlineheadsection
{Frameworks \& Libraries:}
{Spring Boot, Netflix \acr{OSS}, CakePHP, jQuery, Angular, React.}

\inlineheadsection
{Infrastructure \& DevOps:}
{\acr{AWS}, Azure, Google Cloud, Terraform, Docker, Kubernetes, Jenkins, SonarQube, GitLab \acr{CI/CD}, GitHub Actions, NGINX, Apache, Tomcat, Postgres\acr{SQL}, My\acr{SQL}, \acr{SQL}Server.}

\inlineheadsection
{Methodologies:}
{Agile, Scrum, Domain-Driven Design, Clean Architecture, Microservices.}

\inlineheadsection
{Natural languages:}
{Spanish \textit{(native)}, Guarani \textit{(native)}, English \textit{(full professional proficiency)}, German \textit{(limited working proficiency)}, French \textit{(elementary proficiency)}.}

\inlineheadsection
{Certifications:}
{System Thinking --- MIT xPRO \textit{(2024)}, Nanotechnology for Health --- imec \textit{(2013)}, Introduction to Artificial Intelligence --- Udacity \textit{(2009)}.}

%------------------------------------------------

\end{document}
